\chapter{定値性}
    \section{正定Hermite対称行列の逆行列も正定}
        \begin{proof}
            \quad\par
            $A\in\complexNumbers^{n\times n}$をHermite対称な正定行列とし，重複を含めた$n$個の固有値を$\lambda_1,\dotsc,\lambda_n > 0$とする。
            $A$はHermite対称だから適当なユニタリ行列$P\in\complexNumbers^{n\times n}$が存在して$A = P\diag{\lambda_1,\dotsc,\lambda_n}\HerConj{P}$と表せるので
            $A^{-1} = P\diag{\lambda_1,\dotsc,\lambda_n}^{-1}\HerConj{P} = P\diag{1/\lambda_1,\dotsc,1/\lambda_n}\HerConj{P}$である。
            よって$A^{-1}$は対称で固有値が全て正なので正定である。
        \end{proof}
    \section{\texorpdfstring{$A$}{A}が半正定で\texorpdfstring{$\vx^\top A\vx=0$}{x\string^⊤Ax = 0}ならば\texorpdfstring{$\vx\in\Ker{A}$}{x∈Ker(A)}}
        \begin{proof}
            \quad\par
            $A$を$n\times n$実対称行列とする。
            $A$を対角化する実直交行列を$P = [\bm{p}_1,\dotsc,\bm{p}_n]$とし，$\bm{p}_i$に対応する固有値を$\lambda_i$とすると，$A$が半正定だから$\lambda_i\geq 0$である。
            $\vx$を$\realNumbers^n$の正規直交基底$\{\bm{p}_1,\dotsc,\bm{p}_n\}$で展開して$\vx = \sum_{i=1}^n c_i\bm{p}_i$とすると$0 = \vx^\top A\vx = \sum_{i=1}^n \lambda_i {c_i}^2$だから正の固有値に対応する$c_i$は全て0である。
            ゆえに$\vx$は$A$の固有値$0$の固有空間すなわち$\Ker{A}$に属する。
        \end{proof}
    \section{Gram行列は半正定}
        \begin{proof}
            \quad\par
            $V$を内積空間とし，$\vv_1,\dotsc,\vv_n\in V$のGram行列を$G=[\inProd{\vv_i}{\vv_j}]_{n\times n}$とする。
            $\vx\in \complexNumbers^n$に対して
            \[ \vx^\top G\vx = \sum_{i=1}^n\sum_{j=1}^n x_i x_j \inProd{\vv_i}{\vv_j} = \sum_{i=1}^n\sum_{j=1}^n \inProd{x_i\vv_i}{x_j\vv_j} = \inProd{\sum_{i=1}^n x_i\vv_i}{\sum_{j=1}^n x_j\vv_j} = \norm{\sum_{i=1}^n x_i\vv_i}_2^2 \geq 0 \]
        \end{proof}
    \section{正定値行列同士の積は正定値とは限らない}
        \begin{proof}
            \quad\par
            (例) : $75^\circ$左回転行列同士の積\\
            (※ : 行列が正定値というのは，飛ばす前のベクトルと飛ばしたあとのベクトルの成す角が$90^\circ$未満であることと同値)
        \end{proof}
    \section{対角成分が全て非負の対称行列が半正定とは限らない}
        \begin{proof}
            \quad\par
            例えば
            $
                A =
                \left[
                    \begin{array}{cc}
                        1 & -2\\
                        -2 & 1\\
                    \end{array}
                \right]
            $
            とすると$x_1=x_2$なる$\vx$に対して$\vx^\top A\vx = -2{x_1}^2 < 0$
        \end{proof}
    \section{非対称行列は固有値が全て正でも正定とは限らない}
        \begin{proof}
            \quad\par
            例えば
            $
                A =
                \begin{bmatrix}
                    1 & -3 \\
                    0 & 1
                \end{bmatrix}
            $
            の固有値は全て$1$であるが
            $
                [1,1]A
                \begin{bmatrix}
                    1 \\
                    1
                \end{bmatrix}
                = -1
            $
        \end{proof}
    \section{\texorpdfstring{$A$}{A}: Hermite対称で正定\texorpdfstring{$\Rightarrow (\HerConj{\vx}A\vx)(\HerConj{\vy}A^{-1}\vy) \geq |\HerConj{\vx}\vy|^2$}{⇒ (x\string^*Ax)(y\string^*A\string^(-1)y) ≧ |x\string^*y|\string^2}}
        \begin{shadebox}
            $A \in \complexNumbers^{n\times n}$はHermite対称で正定であるとする。
            このとき，任意の$\vx, \vy \in \complexNumbers^n$に対して次式が成り立つ。
            \[ (\HerConj{\vx}A\vx)(\HerConj{\vy}A^{-1}\vy) \geq \|\HerConj{\vx}\vy\|_2^2 \]
            等号が成立するのは$\vx$と$\vy$どちらかが$\bm{0}$であるか，そうでなければ$\vy$と$A\vx$が平行であるとき，かつそのときである。
        \end{shadebox}
        \begin{proof}
            \quad\par
            $A$はHermite対称で正定であるから，適当なユニタリ行列で対角化可能であり，固有値が全て正であることに注意すれば，$\sqrt{A},\sqrt{A^{-1}}$を定義できる。
            これを用いると次式を得る。
            \[ (\HerConj{\vx}A\vx)(\HerConj{\vy}A^{-1}\vy) = \norm{\sqrt{A}\vx}_2^2 \norm{\sqrt{A^{-1}}\vy}_2^2 \geq \left|\inProd{\sqrt{A}\vx}{\sqrt{A^{-1}}\vy}\right|^2 = |\HerConj{\vx}\vy|^2 \]
            2つ目の不等号はCauchy-Schwarzの不等式より導かれる。
        \end{proof}